% Sinh boi code/gen_response_tcom.py tu reviews/so-nhan-xet.csv. KHONG sua tay.
% ⛔ NOP DANG SUPPLEMENTARY FILE, khong phai cover letter (yeu cau cua editor).
\documentclass[10pt]{article}
\usepackage[margin=1in]{geometry}
\usepackage{amsmath,amssymb}
\usepackage[colorlinks=true,linkcolor=black,urlcolor=blue]{hyperref}
\usepackage{enumitem}
\setlist{nosep}
% macro dung chung voi main.tex: thieu chung thi LaTeX nem loi ma bo sinh
% van bao thanh cong. Da vap 03/10: 21/21 muc nhung 4 loi Undefined control sequence.
\newcommand{\Rf}{R_f}
\newcommand{\rs}{R_s}
\newcommand{\Pe}{P_{\mathrm{err}}}
\newcommand{\Pc}{P_{\mathrm{corr}}}
\newcommand{\Psift}{P_{\mathrm{sift}}}
\newcommand{\QBER}{\mathrm{QBER}}
\emergencystretch=1em
\newcommand{\rev}[1]{\par\medskip\noindent\textit{#1}\par\smallskip}
\newcommand{\act}[2]{\noindent\textbf{Action.} #1\par\noindent\textbf{Where.} #2\par}
\title{Response to Reviewers\\ \large TCOM-TPS-26-1667 --- First Revision}
\author{\normalsize KAN-Based Adaptive Parameter Control for Multi-User Satellite FSO/QKD Systems}
\date{}
\begin{document}
\maketitle

\section*{Before the point-by-point replies}

We thank the Associate Editor and the four reviewers. This letter answers \textbf{every} numbered point: 5 from Reviewer~1, 3 from Reviewer~2, 5 from Reviewer~3, 8 from Reviewer~4, 21 in total. It is generated from a ledger with one row per comment, so the count can be checked rather than trusted.

\subsection*{Reviewer 1 --- 5 comments}

\paragraph{Comment 1.}
\rev{The motivation for selecting KAN over other established learning models is not sufficiently convincing. Although KAN offers improved parameter efficiency and symbolic interpretability, the performance gains over the regularized MLP baseline appear relatively modest. Explain why KAN is particularly well suited to this optimization problem and UNDER WHAT CONDITIONS its advantages become significant.}
\act{This is the question the revision is built around, and we can now answer the second half of it (\emph{under what conditions does the advantage become significant}) with a measurement rather than an argument. The controller has four outputs. Three of them are \emph{constant} over the operating regime, because the optimal modulation depth saturates at the physical ceiling of the intensity-modulation index: sweeping it on a $501$-point grid up to $0.999$, the optimum sits at the upper boundary in $80\%$ of sampled channel conditions, and widening the physical sampling range makes that worse, not better. On a constant target no model can outperform any other, and a nine-parameter linear map is exactly as accurate as a network. The advantage of a learned model is therefore confined to the one output that genuinely varies, and there it is clear: per-user accuracy $0.010 \pm 0.001$ for KAN against $0.018 \pm 0.022$ for the regularised MLP and $0.083$ for the linear baseline, with $589$ parameters against $1386$. We have added this characterisation to the results and stated the saturation explicitly in the limitations, because a reader deciding whether to deploy a learned controller needs to know which axis is worth learning.}{Sec.~VI, Table~V; Sec.~VII}

\paragraph{Comment 2.}
\rev{The optimization framework is primarily supported by empirical evidence. The decomposition is numerically validated on small cases, but the ABSENCE OF THEORETICAL CONVERGENCE OR OPTIMALITY GUARANTEES limits the technical rigor. A more comprehensive theoretical discussion or stronger justification of the decomposition would strengthen the contribution.}
\act{Strengthened where a proof exists, and sharpened where one does not. The two halves of the solver had been carrying the same word, \emph{exact}, while having very different status, and we have separated them. Section~IV now gives the one-line argument that the outer/inner split is lossless: with $\boldsymbol{\beta}$ ranging over a compact box and the objective continuous on it, partial maximization gives $\max \Rf = \max_{(\mu,\chi)}[\max_{\boldsymbol{\beta}} \Rf]$ exactly. No relaxation and no assumption beyond compactness is involved, so this part needs no empirical support at all. What that identity does not give is that our inner routine attains the inner maximum. We now name the exact step that lacks a guarantee: the per-user one-dimensional maximizations are exhaustive on a grid plus refinement and so attain their own optimum, while the \emph{mean-field coupling between users} has no contraction proof. For that step we report evidence and label it as evidence: agreement with a brute-force joint search at $N{=}2$, convergence within six iterations, and the initialization study added for Reviewer~4. We would rather mark that boundary than let one adjective cover both sides of it.}{Sec.~IV}

\paragraph{Comment 3.}
\rev{Expand the experimental evaluation to demonstrate robustness and generalization: additional studies under different atmospheric conditions, satellite geometries, user distributions, or orbital configurations.}
\act{Added as a new subsection, covering all four axes you named. We re-solved the oracle on $25$ fresh clusters under each of sixteen conditions: turbulence (wind $11$ to $41$~m/s in the Hufnagel--Valley profile), serving geometry (zenith bands from $0$--$4^\circ$ to $20$--$35^\circ$), cluster size ($N{=}2$ to $6$), and eavesdropper standoff ($10$ to $120$~m). Two things survive, and the second was not what we expected. \textbf{First}, the modulation-depth saturation is a property of the problem rather than of one regime: the optimum sits at its upper bound in $100\%$ of clusters in every feasible condition, across every turbulence level, zenith band and cluster size. The one exception is a close eavesdropper, where at a $10$--$20$~m standoff the URA constraint binds and only $52\%$ remain at the bound. \textbf{Second}, turbulence \emph{helps} security here, which we report as a design constraint and not as a result in our favour. Two conditions yielded no feasible cluster: the $20$--$35^\circ$ zenith band, for the expected reason of a longer slant path, and the \emph{low}-turbulence case at $11$~m/s. The reason is that a beam splitter is detected only through the loss it imposes on the sifting statistics, so detectability scales with the width of the fading distribution. At a fixed geometry $m_{\mathrm{BSA}}$ rises monotonically with turbulence: $0.0038$ at $11$~m/s, which is below the $0.005$ threshold, then $0.0063$, $0.0089$ and $0.0105$ at $21$, $31$ and $41$~m/s. A clear, still night is therefore the worst case for this criterion, not the best.}{Sec.~VI, robustness subsection}

\paragraph{Comment 4.}
\rev{Although multiple random seeds are reported, FURTHER STATISTICAL ANALYSIS would improve confidence in the reported performance differences, particularly between KAN and the MLP baseline.}
\act{Expanded from five seeds to twenty, and more importantly we now report what the spread means. The earlier spread was dominated by the two software faults described in our reply to Reviewer~2, not by model variability: once they are fixed, the seed standard deviation of the per-user accuracy drops from $0.034$ to $0.001$ and the closed-loop retention becomes stable to three decimals. We report mean and standard deviation over twenty seeds, and we additionally report the number of divergent cells per configuration, because a zero there is what makes the mean meaningful.}{Sec.~VI, Table~V}

\paragraph{Comment 5.}
\rev{A more explicit discussion of its LIMITATIONS, including the dependence on ORACLE-GENERATED TRAINING DATA, the relatively LIMITED TRAINING DATASET, and the potential impact of DISTRIBUTION SHIFTS during real-world deployment.}
\act{The limitations subsection has been rewritten rather than extended, because one of its existing entries turned out to be wrong in a way worth flagging. It attributed the cluster controller's seed variance to the oracle sitting on a sharp four-dimensional feasibility boundary. That was a scientific explanation for what we have since traced to two software faults, and we have replaced it with the mechanism and the measurements. On oracle dependence specifically, which you asked about: the controller is trained by supervision from the decomposed solver, inherits whatever that solver gets wrong, and cannot by construction exceed it. We now state this together with a consequence we would otherwise have been tempted to present as a result: once the faults are fixed, every controller we tested, including a nine-parameter linear map, reaches the oracle rate on every test configuration, so the closed-loop retention metric has saturated and no longer separates architectures. Three further limitations are now stated explicitly: the degeneracy of three of the four targets and what it costs; the scope of the eavesdropper model, which covers two attacks modelled individually; and the sharpness of the beam-splitting detection threshold, which spans about a quarter of a percentage point.}{Sec.~VII, limitations}

\subsection*{Reviewer 2 --- 3 comments}

\paragraph{Comment 1.}
\rev{Section III.B, point 3) Inter-user secrecy: the authors do not really justify CAPPING TAU AT 0.1. This is an important bound to information available to attackers, and the section would benefit from an explanation of the CONSEQUENCES of such a cap.}
\act{Justified by measurement, and the question exposed an error we are glad to correct. While preparing the justification we found that the manuscript stated $\tau{=}0.1$ while every run actually used $\tau{=}0.05$; no result in the paper was produced at $0.1$. We have corrected the text to the value that was executed. We then swept $\tau$ over the oracle solutions of all $250$ clusters: the feasible-pair count and the aggregate key rate are \emph{identical} for every $\tau \ge 0.05$, and the constraint only becomes active below $\tau \approx 0.04$, where the feasible count falls from $766$ to $608$. The cap is therefore an inactive guard in the studied regime rather than a tuned parameter, and we now say so instead of asserting a value. The correction changes no number elsewhere in the paper, because $0.05$ is what produced them.}{Sec.~III, inter-user secrecy constraint}

\paragraph{Comment 2.}
\rev{Section IV.F, Table IV: the standard deviations for KAN+resid and MLP+resid KeyRet are OF THE SAME MAGNITUDE AS THE RESULT ITSELF. This is a critical benchmark and the authors' method is OVERCOME BY MLP here (by a low margin), but both are within each other's standard deviations. The high standard deviation is shortly addressed but NOT SATISFACTORILY. If the source is seed-to-seed variance, there might be a BETTER PARAMETER th\dots}
\act{You identified a real defect, and chasing it down changed the paper. The standard deviations were of the same magnitude as the values because two software faults, not two modelling choices, were corrupting the closed-loop evaluation. \textbf{First}, the feature standardiser divided by a per-feature standard deviation guarded only against exact zero. One controller input is constant in this regime, with a standard deviation of $7.99\times10^{-15}$ rather than zero, so the guard was bypassed and floating-point noise was amplified by a factor of $1.25\times10^{14}$; predicted thresholds then moved by $0.35$ in response to last-bit input differences. \textbf{Second}, three of the four controller targets are constant in the studied regime, so the residual the network was asked to fit had a standard deviation of $1.74\times10^{-16}$; the quasi-Newton optimiser diverged on $20/20$ seeds, producing $6320$ non-finite cells that a silent fallback replaced with the training median. A diverged model was therefore being reported as a measurement. Both are fixed: the standardiser now uses a relative tolerance, and constant targets are emitted exactly rather than learned. Divergent cells fall from $6320$ to $0$, the seed standard deviation of the per-user accuracy falls from $0.034$ to $0.001$, and the reported counts become reproducible across machines. We also report the divergence count alongside every row, so a diverged run can never again be read as a measurement.}{Sec.~VI, Table~V; artifact \texttt{scripts/train\_cluster.py}}

\paragraph{Comment 3.}
\rev{KAN and MLP have close results in several benchmarks. The paper would benefit from a MORE DETAILED DESCRIPTION OF THESE DIFFERENCES and WHERE EACH APPROACH IS BETTER SUITED for the architecture.}
\act{Described, and the honest description is not the one we would have written before the revision. With the two software faults fixed, KAN and the MLP are \emph{indistinguishable} on the closed-loop metric: both reach the oracle rate on every test configuration, as does the linear baseline, so that metric has saturated and can no longer separate the models. We say so rather than quoting it as a win. Where they do separate is accuracy on the single non-degenerate output and the stability of that accuracy. KAN gives $0.010 \pm 0.001$; the MLP gives $0.018 \pm 0.022$, where the across-seed standard deviation \emph{exceeds} the gap between the two means. The difference that matters operationally is therefore not that KAN is on average better but that the MLP is erratic across initialisations while KAN is not, at $43\%$ of the parameter count. Our reading: the MLP is the better choice when the target has broad support and retraining is cheap; the KAN is the better choice here, where the learnable signal is confined to one output, the parameter budget is on-board, and a single bad initialisation cannot be detected in flight.}{Sec.~VI, Table~V}

\subsection*{Reviewer 3 --- 5 comments}

\paragraph{Comment 1.}
\rev{Whether the two considered eavesdropping strategies (URA and BSA) provide a COMPREHENSIVE analysis of Eve's condition. Further, a MORE DETAILED DESCRIPTION OF BSA is needed: Eve's channel realization details, transmittance capacity of the beam-splitter, etc.}
\act{Two separate answers, and the first one is no. \textbf{On comprehensiveness:} URA and BSA are \emph{not} a comprehensive treatment of Eve, and the manuscript should not have left that implicit. They are the two attacks for which this physical layer admits a closed-form detectability criterion, and they are modelled \emph{individually}, not jointly. Not covered: photon-number-splitting on the weak-coherent source, Trojan-horse probing of the modulator, detector blinding or time-shift attacks on the dual-threshold receiver, side channels in classical post-processing, and any collective or coherent attack. We now say this explicitly rather than by omission, and list those classes as future work. Section~I also narrows the security language accordingly, as Reviewer~4 asked. \textbf{On the BSA detail:} we have added what the criterion actually depends on. The splitter is passive and draws a fixed power fraction $\eta_{\mathrm{leak}}$ at the relay; detectability enters only through the loss it imposes on the legitimate sifting probability, $m_{\mathrm{BSA}}(\beta)=1-\Psift(\gamma(1-\eta_{\mathrm{leak}}))/\Psift(\gamma)$, so Eve's own receiver is \emph{not} modelled, only the trace she leaves. We also swept the splitter ratio, which the manuscript had fixed without justification: the fraction of pairs for which detection can be certified falls from $100\%$ at $\eta_{\mathrm{leak}}\ge1.4\%$ to $73\%$ at $1.25\%$ and to zero at $1.15\%$, so the whole transition occupies about $0.25$ percentage points. The design therefore certifies detection of a splitter drawing $1.4\%$ or more; against a stealthier one the constraint is not met and no key is emitted. That is the conservative behaviour, but it is also a real limit, and it is now stated.}{Sec.~III, URA and BSA subsections; Sec.~VII}

\paragraph{Comment 2.}
\rev{This work considers modulation-signalling based QKD, a very active research topic, but the LITERATURE SURVEY LOOKS QUITE STRAIGHTFORWARD, providing less detailed information about recent relevant research. Add recent relevant references.}
\act{Expanded, and we tried to make the addition do work rather than lengthen a list. The survey now covers the modulation-signalling branch specifically and says, for each entry, what it settles and what it leaves open for us. Subcarrier-wave QKD has been analysed numerically for free-space ship-to-port links \cite{goncharov2025scw} and routed through a deployed metropolitan transport network \cite{tarabrina2023scw}, which puts the signalling format outside the laboratory but in both cases on a \emph{single} link rather than a shared relay. Direct intensity and phase modulation has been demonstrated as a QKD transmitter \cite{liu2026intensity}, which is the hardware assumption our modulation depth parameterizes. The free-space continuous-variable line is surveyed in \cite{kargina2026cvqkd}; it shares our reliance on classical coherent components but replaces the dual-threshold decision with homodyne detection, so the parameter that dominates our problem has no counterpart there. We also added two review articles at Reviewer~4's suggestion and a recent learning-based controller at their request. Every entry was verified against Crossref before being cited.}{Sec.~I, paragraphs on the modulation-signalling branch}

\paragraph{Comment 3.}
\rev{FIGURE 4 NEEDS TO BE IMPROVED. The legends are confusing and appear at a wrong place. Correct font size and correct legend placement can improve the presentation a lot.}
\act{Redrawn. The single figure-level legend sat above the panel titles and described the bars of the lower row while being placed over the upper row, whose curves use different colours and were not in the legend at all. The upper row now carries its own legend inside the panel, the bar legend sits directly beneath the lower row, and font sizes were raised. We also corrected something the review did not raise: the inset box reported a step count next to a label that named the key-rate gain, mixing two different quantities; the inset now names each explicitly.}{Fig.~4 and its caption}

\paragraph{Comment 4.}
\rev{The authors compare KAN, MLP and a linear baseline; a discussion is needed on HOW JUSTIFIED it is to use and compare these models AS COMPARED TO OTHER RELEVANT SCHEMES.}
\act{Justified, and the revision strengthened the case for keeping the linear baseline rather than weakening it. Three of the four controller outputs are constant in this regime, and on those a nine-parameter linear map attains the same accuracy as either network; without that baseline in the table a reader could not see that, and would credit the networks for accuracy that requires no network at all. The linear model also reaches the oracle rate on every test configuration, so it is not a straw man. We kept the regularised MLP as the like-for-like learned competitor and report its parameter count alongside, since parameter efficiency is one of the axes under comparison. We did not add further learned baselines: tuning a third architecture on our scenario would make the comparison depend on how well we tuned someone else's model, and we would rather state that limit than hide it behind an under-tuned row.}{Sec.~VI, Table~V}

\paragraph{Comment 5.}
\rev{A few minor errors: reference [9] has the AUTHOR'S NAME MISSING. The terms 2.10x, 3.6x in the abstract should use a word like TIMES. There are many TYPO ERRORS AND GRAMMATICAL MISTAKES that need attention.}
\act{All three fixed. (i) Reference [9] showed \texttt{------} because the bibliography style abbreviates a repeated author list; we disabled that behaviour, and the entry now prints the full author list. (ii) The abstract no longer uses the $\times$ symbol: the gains are written as \emph{times}. (iii) We made a consistency pass over the manuscript and corrected a duplicated word in Section~I, two abbreviations that were re-expanded after their first definition, and an inconsistent rendering of the name Kolmogorov--Arnold.}{References; Abstract; Sec.~I; Sec.~II}

\subsection*{Reviewer 4 --- 8 comments}

\paragraph{Comment 1.}
\rev{Towards the end of the first paragraph of the introduction, include the following review articles to make the literature review on satellite quantum communication, particularly satellite QKD, more comprehensive: RevModPhys.92.025002 and RevModPhys.94.035001.}
\act{Both reviews are now cited at the end of the first paragraph, and we say what each contributes rather than appending them to a list. We note that \cite{xu2020rmp} surveys how device imperfections, not protocol design, set the security of practical QKD, which is precisely the regime our parameter-control problem lives in; and that \cite{lu2022micius} reviews the Micius programme, still the only platform on which satellite QKD has been demonstrated end to end, and the source of the link budgets that make a multi-user relay architecture worth optimizing. We verified both entries against Crossref.}{Sec.~I, end of first paragraph}

\paragraph{Comment 2.}
\rev{The manuscript CLAIMS INFORMATION-THEORETIC SECURITY yet adopts simplified wiretap-like metrics based on Eve's symbol error probability and mutual information. A very small mutual information does not necessarily imply operational security, and security proven only through accessible information is generally INSUFFICIENT FOR MODERN QKD SECURITY. Estimate a bound on the FAILURE PROBABILITY of the cryptographic system like in CO\dots}
\act{We agree, and we have narrowed the claim rather than defended it. The phrase \emph{information-theoretic security} no longer appears as a property we establish. Section~I now states the scope precisely: the secret-key rate used throughout is the asymptotic Csisz\'ar--K\"orner fraction $\max\{0, H_2(P_e)-H_2(\mathrm{QBER})\}$ per sifted bit, evaluated against the two attack models we model; it carries no finite-key correction, no composable-security statement, and no claim against general coherent attacks. The guarantee reported in the prior work is attributed to that work, under its own assumptions.}{Sec.~I, paragraph 2}

\paragraph{Comment 3.}
\rev{The stated convergence analysis being EMPIRICAL IS WEAK. The manuscript states the mean-field iteration converges in 6 iterations and that a unique fixed point was ONLY OBSERVED EXPERIMENTALLY. For this journal a stronger theoretical support is expected: MONOTONICITY ARGUMENTS, FIXED-POINT EXISTENCE CONDITIONS, LOCAL CONVERGENCE ANALYSIS, and SENSITIVITY TO INITIALIZATION, wherever possible.}
\act{You were right that the convergence evidence was weak, and testing it properly overturned one of our statements. We report what we can prove and what we measured, separately. \textbf{Proved:} the outer/inner split is a lossless partial maximization, which needs only compactness of the threshold box and continuity of the objective; the argument is now given in Section~IV. The per-user one-dimensional subproblems are solved by exhaustive grid plus refinement and so attain their own optimum. \textbf{Measured, and it contradicts our earlier claim:} we had written that the mean-field fixed point was unique on every cluster tested. Restarting the inner iteration from one common-threshold value at a time, \emph{all 12 of 12} tested cluster states reach \emph{different} fixed points depending on the start, with the Alice-side threshold differing by up to $1.31$. The uniqueness we observed was an artefact of the solver already running four starts and keeping the best. We have retracted the uniqueness statement, and we now say that multi-start is a requirement rather than a convenience. \textbf{Not provided:} a contraction proof or local convergence rate for the mean-field coupling. We would rather state that gap and give the counterexample count than offer a monotonicity argument we cannot establish for the coupled map.}{Sec.~IV; artifact \texttt{scripts/do\_trien\_khai.py}}

\paragraph{Comment 4.}
\rev{In the complexity discussion the manuscript reports about 4 s per outer sweep node and roughly 200 s per analyzed pass. Add a discussion on PRACTICAL DEPLOYMENT TIMING REQUIREMENTS, HOW OFTEN re-optimization must occur, and IF OR WHEN RE-TRAINING CAN BE AVOIDED.}
\act{Added, and the measurement argues against one motivation we might have claimed. On eight independent passes we recomputed the optimum at every step and then asked what it costs to hold parameters for longer. Refreshing every $50$~s is the reference; holding for $100$~s retains $88.9 \pm 3.5\%$ of the key, $150$~s retains $84.4 \pm 2.0\%$, $250$~s retains $68.3 \pm 9.4\%$, and beyond about $600$~s the link is mostly lost. Re-optimization every one to two steps is therefore the operating requirement. The oracle itself costs $7.57 \pm 2.15$~s per step on the reported server, which is well inside that budget. We therefore do \emph{not} claim the learned controller is needed for latency on this hardware, and we have removed any suggestion of it. Its case rests on operating without the full physics model and channel state at the node, and on a bounded, deterministic inference cost. One further observation we did not expect: the retention curve is not monotone in the refresh interval, because what matters is where the refresh instants fall relative to the geometry of the pass, not how many of them there are. A fixed-clock refresh policy is therefore the wrong design; refresh should be triggered by change in the channel state. We say so rather than quoting a single cadence number.}{Sec.~VI; Sec.~VII}

\paragraph{Comment 5.}
\rev{If possible, add a COMPARISON WITH REINFORCEMENT LEARNING-BASED adaptive QKD approaches, for instance arxiv.org/html/2603.04192v2.}
\act{We have added and discussed the suggested work, but we have not added it as a numerical baseline, and we would rather give the reason than appear to dodge. Section~I now contrasts it on three axes. Its physical layer is photon-counting discrete-variable QKD with decoy intensities as the tuned parameters; ours is intensity-modulated direct detection with a modulation depth and two detection thresholds, so the parameter spaces do not overlap and a head-to-head number would compare protocols rather than controllers. It is single-link, so the coupling that makes our problem hard, one global modulation depth shared across the cluster, does not arise. And it learns online by reinforcement whereas we learn offline by supervision from a decomposed oracle. A reinforcement-learning controller for the multi-user FSO setting is a comparison worth making, and we state it as future work: constructing a fair one is a study in itself, and a weak re-implementation would be worse evidence than none.}{Sec.~I, paragraph before the contributions; Sec.~VII}

\paragraph{Comment 6.}
\rev{The NOTATIONS BECOME CUMBERSOME in Sections III to V. A NOTATION TABLE would greatly improve readability.}
\act{Added. Table~I collects every symbol used in Sections~II--IV, grouped into decision variables, geometry and channel, detection and key rate, and solver and controller. Each entry was checked against the body text, so the table cannot list a symbol the paper does not use.}{Table~I, start of Sec.~II}

\paragraph{Comment 7.}
\rev{In the results, many reported gains are based on averages over ONLY FIVE RANDOM SEEDS. Considering the large observed variances, including CONFIDENCE INTERVALS OR HYPOTHESIS TESTS would be useful.}
\act{Agreed, and we have rerun everything at twenty seeds. We also want to be explicit about something the extra seeds revealed rather than hid: at five seeds the closed-loop numbers were not reproducible across machines at all, because a standardisation fault made them sensitive to the fifteenth decimal digit of a prediction. We traced that causally, by perturbing predictions by a controlled relative amount and watching the feasible count move, and we fixed the cause rather than averaging over it. The twenty-seed table is therefore not simply a larger sample of the old experiment; it is the first run whose numbers are machine-independent.}{Sec.~VI, Table~V}

\paragraph{Comment 8.}
\rev{The cluster controller is trained on approximately 250 CLUSTER SAMPLES. The reported standard deviations 0.48 for KAN and 0.40 for MLP indicate SUBSTANTIAL INSTABILITIES, raising concerns on INSUFFICIENCY OF THE DATASET SIZE. Either EXPAND THE DATASET or comment on the LIMITS OF GENERALIZABILITY, ROBUSTNESS AND DEPLOYMENT READINESS.}
\act{The sample size is unchanged at $250$ clusters, and we now state its consequences rather than leaving them implicit. With the two faults fixed, the across-seed standard deviations are $0.001$ for the per-user accuracy and $\le 0.004$ for closed-loop retention, so sample size is no longer the limiting source of uncertainty for those quantities. What $250$ clusters does limit is coverage of the geometry space, and we say so in the limitations: the clusters are drawn from a single sampling regime, and we make no claim about behaviour outside it.}{Sec.~VI; Sec.~VII, limitations}

\end{document}
